\documentclass[a4paper,twoside]{article}

\usepackage{morewrites}

\usepackage[svgnames,dvipsnames,x11names]{xcolor}
\usepackage{graphicx}
\usepackage[english]{babel}
\usepackage[colorlinks=true, linkcolor=NavyBlue, urlcolor=NavyBlue, citecolor=NavyBlue]{hyperref}
\usepackage[a4paper]{geometry}
\usepackage{ts-fonts}
\usepackage{tcolorbox}
\usepackage{fvextra}
\usepackage{amsmath}
\usepackage{float}
\usepackage{seqsplit}
\usepackage{ts-code}

\usepackage{lastpage}
\usepackage{refcount}
\makeatletter
\def\ps@plain{
\let\@oddhead\@empty
\let\@evenhead\@empty
\def\@oddfoot{
\hfil
\ifnum\getpagerefnumber{LastPage}>1\relax
\thepage
\fi
\hfil}
\let\@evenfoot\@oddfoot
}
\makeatother

\title{Mermaid Configuration}
\author{}\date{}
\begin{document}
\maketitle
TeXSmith will automatically pick up a \tscodeinline{mermaid-\allowbreak{}config.json} located at the root of a template (next to \tscodeinline{manifest.toml}). The diagrams module passes this config to Mermaid for all diagrams rendered with that template.
\section{Using a Built-in Template}
The built-in \tscodeinline{article} template ships with a \tscodeinline{mermaid-\allowbreak{}config.json}. To inspect or override it:
\begin{tscode}[lang=bash, engine=pygments]
\begin{Verbatim}[commandchars=\\\{\},breaklines, breakanywhere, commandchars=\\\{\}]
texsmith\PY{+w}{ }\PYZhy{}\PYZhy{}template\PYZhy{}info\PY{+w}{ }\PYZhy{}\PYZhy{}template\PY{+w}{ }article
texsmith\PY{+w}{ }templates\PY{+w}{  }\PY{c+c1}{\PYZsh{} list all discoverable templates}
\end{Verbatim}
\end{tscode}
To customize, copy the file, adjust options, and point to your modified template directory:
\begin{tscode}[lang=bash, engine=pygments]
\begin{Verbatim}[commandchars=\\\{\},breaklines, breakanywhere, commandchars=\\\{\}]
cp\PY{+w}{ }\PYZhy{}r\PY{+w}{ }\PY{k}{\PYZdl{}(}python\PY{+w}{ }\PYZhy{}\PY{+w}{ }\PY{l+s}{\PYZlt{}\PYZlt{}\PYZsq{}PY\PYZsq{}\PYZbs{}nfrom texsmith.core.templates import load\PYZus{}template\PYZbs{}nfrom pathlib import Path\PYZbs{}nt = load\PYZus{}template(\PYZsq{}article\PYZsq{})\PYZbs{}nprint(t.root)\PYZbs{}nPY}\PY{k}{)}\PY{+w}{ }./templates/article\PY{l+s+se}{\PYZbs{}n}\PY{c+c1}{\PYZsh{} edit ./templates/article/mermaid\PYZhy{}config.json\PYZbs{}ntexsmith doc.md \PYZhy{}\PYZhy{}template ./templates/article}
\end{Verbatim}
\end{tscode}
\section{Adding Mermaid Config to a Custom Template}
\begin{enumerate}
\item Place \tscodeinline{mermaid-\allowbreak{}config.json} at the template root (same level as \tscodeinline{manifest.toml}).
\item TeXSmith will expose the path via \tscodeinline{template.extras["mermaid\_config"]} so the renderer can pass it to Mermaid.
\item No manifest changes are required; the presence of the file is enough.
\end{enumerate}
Typical options include theme, font, backgroundColor, and securityLevel. See \url{https://mermaid.js.org/config/theming.html} for full reference.
\end{document}
